\begin{workedexample}[Worked Example: Cable loss budget]
A $10\ \text{m}$ coax run is rated $0.2\ \text{dB/m}$ at $1\ \text{GHz}$,
so the total attenuation is $0.2\times10 = 2\ \text{dB}$: a $+10\ \text{dBm}$
signal arrives at $+8\ \text{dBm}$ ($6.3\ \text{mW}$ of the original $10\ \text{mW}$,
about $63\%$). With velocity factor $0.66$, the one-way delay over the run is
\[ t = \frac{\ell}{\text{VF}\cdot c} = \frac{10}{0.66\times3\times10^{8}} = 50.5\ \text{ns}. \]
\end{workedexample}

\begin{keytakeaways}
\begin{itemize}[leftmargin=1.4em,itemsep=2pt]
\item Total loss (dB) $=$ per-metre loss $\times$ length; dB losses add.
\item Conductor loss grows as $\sqrt{f}$ (skin effect); dielectric loss grows roughly linearly with $f$.
\item Delay $t = \ell/(\text{VF}\cdot c)$; a lower velocity factor means more delay.
\end{itemize}
\end{keytakeaways}

\begin{exercises}
\item A cable loses $0.15\ \text{dB/m}$. What is the total loss over $20\ \text{m}$?
\item What fraction of the power survives $3\ \text{dB}$ of loss?
\item What is the delay through $5\ \text{m}$ of coax with velocity factor $0.66$?
\end{exercises}

\furtherreading{Pozar~\cite{pozar} (ch.~2); Ludwig and Bogdanov~\cite{ludwig}.}
